\section{条件概率与乘法法则}
\label{sec:05-Probability/02-Conditional-Probability/01}

你从一副标准扑克牌里抽一张。抽到 A 的概率是 \(4/52=1/13\approx0.0769\)。

这时候旁边有人瞄了一眼你的牌，说了一句：``是黑桃。''

你的概率立刻变了。原来的五十二种可能不再都站得住——红桃、方块、梅花全部出局，只剩下十三张黑桃。而黑桃里恰好有一张 A。所以现在抽到 A 的概率是 \(1/13\approx0.0769\)。

等一下。\(4/52\) 和 \(1/13\) 不都是 \(0.0769\) 吗？这次恰好相等。但换个情况：你抽到的如果是人头牌（J、Q、K），情况就不一样了。

原先抽到人头牌的概率是 \(12/52=3/13\approx0.231\)。被告知是黑桃后，黑桃中的人头牌张数是多少？三张——J、Q、K。所以条件概率是 \(3/13\approx0.231\)。又相等？

再试一个：抽到的牌点数大于等于 \(10\)（即 \(10\)、J、Q、K、A）。原先有 \(5\times4=20\) 张，概率 \(20/52=5/13\approx0.385\)。黑桃中有 \(5\) 张（\(10\)、J、Q、K、A），条件概率 \(5/13\approx0.385\)。

你发现了一个规律：如果只看``是黑桃''这个条件，似乎很多概率都和原来一样？不，这是因为扑克牌的四个花色在点数分布上是完全对称的——每个花色里各点数恰好一张。对称意味着条件不影响相对比例。换个不对称的例子立刻现原形。

\subsection{信息缩小了你的世界}

袋子里有三只红球、两只蓝球。闭眼摸一只，摸到红球的概率是 \(3/5=0.6\)。

第一次摸完不放回，你把手伸进去准备摸第二只。现在袋里的球变了——你第一次摸到的是什么你心里有数。如果第一次摸走了红球，袋中剩两红两蓝，第二次摸到红球的概率是 \(2/4=0.5\)，不是原来的 \(0.6\)。如果第一次摸走了蓝球，袋中剩三红一蓝，第二次摸到红球的概率是 \(3/4=0.75\)。

你没法用一个单一数字回答``第二次摸到红球的概率''，因为答案取决于第一次发生了什么。条件概率做的就是这件事：\textbf{在被告知某个事件已经发生的条件下，重新计算另一个事件的概率}。

\subsection{定义：限制、再缩放}

若 \(P(B)>0\)，在事件 \(B\) 已发生的条件下 \(A\) 的概率定义为

\[
P(A\mid B)=\frac{P(A\cap B)}{P(B)}.
\]

拆开这个公式的两个动作：

\begin{enumerate}
\tightlist
\item 分子 \(P(A\cap B)\)：只保留 \(A\) 和 \(B\) 同时成立的那些结果。\(B\) 之外的所有结果被当场删除。
\item 分母 \(P(B)\)：在 \(B\) 的范围内重新做规范化，让条件事件 \(B\) 自身的概率回到 \(1\)。
\end{enumerate}

用摸球的数字感受：抽两球不放回，第一只是红球的概率 \(P(R_1)=3/5=0.6\)。第一只红的条件下第二只红的概率是 \(P(R_2\mid R_1)=(3/5\times2/4)/(3/5)=2/4=0.5\)。分母 \(3/5\) 约掉之后，剩下的正是第一次摸走红球后残局里红球的比例。

你可能会想：为什么不直接定义 \(P(A\mid B)=P(A\cap B)\)？因为那样的话，所有条件概率加起来不等于 \(1\)——\(B\) 内的各子事件概率之和还是 \(P(B)\)，没有回到全集概率。除以 \(P(B)\) 就是强行把 \(B\) 变成新的``全集''。

用数字检验：\(B=\)第一次红，\(B\) 内的子事件有 \(R_2\) 和 \(B_2\)（第二只是蓝球）。如果把条件概率定义成 \(P(A\cap B)\) 而不除分母，那么 \(P(R_2\cap R_1)+P(B_2\cap R_1)=(3/5\times2/4)+(3/5\times2/4)=3/10+3/10=0.6\)，不等于 \(1\)。除以 \(P(R_1)=0.6\) 之后，\(0.5+0.5=1\)，新全集恢复概率为一。

\subsection{别让分子骗了你}

掷两枚公平骰子。已知点数和为偶数，求第一枚为偶数的概率。

让你先凭直觉猜一个数。点数和为偶数，意味着两枚同奇偶。如果第一枚是偶数，第二枚也得是偶数；第一枚是奇数，第二枚也得是奇数。似乎各半？试试看。

令 \(A=\)第一枚为偶数，\(B=\)点数和为偶数。在三十六个等可能的有序结果中，和是偶数的有多少？两个都是偶数有 \(3\times3=9\) 种，两个都是奇数有 \(3\times3=9\) 种，共 \(18\) 种。所以 \(P(B)=18/36=1/2\)。

第一枚是偶数且和为偶数，意味着第二枚也是偶数：\(3\times3=9\) 种。所以 \(P(A\cap B)=9/36=1/4\)。

\[
P(A\mid B)=\frac{P(A\cap B)}{P(B)}=\frac{1/4}{1/2}=\frac12.
\]

和直觉一致。但很多人第一次看到这个公式时，会直接把分子 \(1/4\) 当成答案——因为他们忘了分母。\(1/4\) 是 \(P(A\cap B)\)，即``第一枚偶数且和为偶数''的概率；而你要的是``在已知和为偶数的前提下，第一枚偶数的概率''。前者是两个条件同时成立的原始概率，后者是在一个条件已经给出的情况下另一个的概率。它们是两个不同的问题。

一个快速检查的方法：算完之后把条件反过来算一遍。已知第一枚为偶数，和为偶数的概率：\(P(B\mid A)=(1/4)/(1/2)=1/2\)。在你这个对称的骰子模型里，两个方向恰好相等，但这不是一般规则。下一个例子会让你看到两个方向的天差地别。

\subsection{乘法法则：把定义翻个面}

条件概率的定义移项：

\[
P(A\cap B)=P(B)P(A\mid B)=P(A)P(B\mid A).
\]

看起来只是一个代数重排，但它对应一种非常自然的计算模式：分步实验。先算第一步的概率，再算第一步发生之后第二步的概率。

袋中三红二蓝，不放回抽两只。两只都是红球的概率：

\[
P(R_1\cap R_2)=P(R_1)P(R_2\mid R_1)=\frac35\cdot\frac24=\frac3{10}=0.3.
\]

第二个因子 \(2/4\)，不是 \(3/5\)。因为你已经拿走了一只红球，残局变了。条件概率强迫你在每次获取新信息后更新残局状态，而不是拿着初始概率一路乘到底。

你可能会想：不放回抽两只红球的概率，能不能直接用组合数？\(\binom{3}{2}/\binom{5}{2}=3/10=0.3\)。数字一样。两种算法对应两种思考习惯：组合数一次性考虑两个位置，乘法法则顺着时间线依次推进。有时候一种比另一种更自然——取决于你怎么想象那个实验过程。关键不是哪种算法更快，而是确保你的概率分配过程与实际物理过程一致。

如果你错误地写成 \((3/5)\times(3/5)=9/25=0.36\)，你实际上假设了第一次摸完把球放回去并重新洗匀。这不是这道题的情况。数字 \(0.3\) 和 \(0.36\) 之间的差距，来自你有没有尊重``不放回''对残局的改变。

\subsection{链式法则：长过程的步步更新}

三个事件：

\[
P(A\cap B\cap C)=P(A)P(B\mid A)P(C\mid A\cap B).
\]

每一项都带着截止到上一步的全部信息。随着链条变长，条件部分越来越长，但每一步的更新逻辑一样：把前面已经发生的全部当作已知，算下一步的条件概率。

一般形式：

\[
P\left(\bigcap_{i=1}^{n}A_i\right)
=P(A_1)\prod_{i=2}^{n}
P\left(A_i\mid A_1\cap\cdots\cap A_{i-1}\right),
\]

只要中间每一步的条件事件概率都不为零。

这条公式什么都不假设——不需要独立，不需要同分布，不需要对称。它对任何一串事件都成立，因为每一步只是条件概率定义的重复使用。独立性会在后面的章节中让某些条件概率简化成无条件概率，但链式法则本身比独立性基础得多。

用数字走一个稍长的链条。一副牌不放回抽五张，前五张都不是 A 的概率：

\[
\frac{48}{52}\cdot\frac{47}{51}\cdot\frac{46}{50}\cdot\frac{45}{49}\cdot\frac{44}{48}.
\]

逐项算：第一张不是 A，52 张中 48 张合格；第二张，残局 51 张中 47 张合格；依此类推。乘起来大约 \(0.659\)。如果你直接数组合数——五张全部来自 48 张非 A 的组合数除以全部五张组合数——即 \(\binom{48}{5}/\binom{52}{5}\)，展开约分之后正好是上面五个因子的乘积。两个视角等价：组合数是一次性空间切分，链式法则是沿着时间线推进，最后落在同一个数字上。

你可能会想：分子和分母最后一项的 \(48\) 不是约掉了吗？没错，\(\frac{48}{52}\times\frac{44}{48}=\frac{44}{52}\)，但中间那几项不能这样直接约。链式法则的真正力量在更复杂的问题中才显现——当每一步的条件概率取决于之前的多步历史，并且不能简化为简单计数时。

\subsection{条件概率不是因果}

\[
P(A\mid B)\neq P(B\mid A)
\]

绝大多数情况下。就算两个数值都很大，也不能推导出因果关系。

\(P(\text{道路湿}\mid\text{下雨})\) 可能高达 \(0.95\)。但 \(P(\text{下雨}\mid\text{道路湿})\) 可能只有 \(0.3\)，因为道路湿的另一个常见原因是洒水车，而洒水车的出现频率远高于下雨。条件概率描述的是信息关联——``知道一件事后另一件事的概率怎么变''——而不是``一件事导致另一件事''。

因果判断需要额外的结构：时间箭头（原因在前）、干预（如果我人为弄湿道路，下雨的概率会不会变？）、以及排除混杂变量。条件概率不做这些事。它只回答一个狭窄而精确的问题：已知 \(B\) 发生时，\(A\) 有多少可能性。把``已知''偷换成``因为''，是概率推理中最常见也最危险的滑移。

你可以这样记住这条边界：条件概率的竖线 \(A\mid B\) 读作``given \(B\)''，不是``because of \(B\)''。竖线是一个信息操作，不是一个因果操作。

\subsection{零概率条件的边界}

条件概率的定义要求 \(P(B)>0\)。如果 \(P(B)=0\)，分母为零，公式不能用。

但概率论和统计学里充满了``给定 \(X=x\)''的表述，而连续随机变量取任意特定值的概率恰好是零。在 \([0,1]\) 上均匀随机选一点，给定 \(X=0.5\) 的概率怎么算？

直接套定义不行。应对方案有几种：对于有联合密度的连续变量，用条件密度 \(f_{Y\mid X}(y\mid x)=f_{X,Y}(x,y)/f_X(x)\)，不直接除概率而除密度；对于更一般的情况，用正则条件概率或通过极限来定义。具体计算在\hyperref[sec:05-Probability/05-Joint-Distributions/01]{``联合、边缘与条件分布''}中展开。

初学时只需带两个意识走：(1) 条件概率定义要求分母非零，但实践中广泛使用的``条件化''操作往往越过了这个边界；(2) 那些越界的用法在连续情形下通常有严格的补救方案，只是补救比定义多了一步，不能直接照搬比值公式。

\subsection{下一站：反向推断}

条件概率竖线左右不能交换，但你最常碰到的恰恰是``需要反向''的情形。你知道 \(P(\text{阳性}\mid\text{患病})\)——检测方法的技术指标，但你想知道的是 \(P(\text{患病}\mid\text{阳性})\)——拿到阳性报告后，你应该多紧张。这两个方向之间隔着一条公式：Bayes 公式。它用全概率公式先算出各种原因产生该结果的总概率，再重新分配解释权重。

下一节会用一个具体的疾病检测案例带你完整走一遍这个反向推断的链条。
